\capitulo{6}{Trabajos relacionados}

\section{Introducción}

Existen numerosas herramientas para crear diagramas relacionados con bases de datos. Sin embargo, no todas persiguen el mismo objetivo ni utilizan la misma notación. Algunas están orientadas al dibujo general de diagramas, otras al modelado profesional de bases de datos y otras al apoyo docente en el aprendizaje del modelo Entidad-Relación.

La comparación realizada en este capítulo no pretende establecer una clasificación absoluta de herramientas, sino situar UBU E-R App frente a alternativas que pueden utilizarse con fines similares. Para ello se consideran criterios como la disponibilidad gratuita o de pago, la facilidad de acceso, la cercanía a la notación Entidad-Relación académica, la validación del modelo, la transformación relacional, la generación SQL y la experiencia de uso.

\section{Herramientas de diagramación general}

diagrams.net, conocido también como draw.io, es una herramienta general de diagramación que permite crear distintos tipos de diagramas, entre ellos diagramas relacionados con software, redes, UML o bases de datos~\cite{drawio}. Su principal ventaja es que ofrece una interfaz madura, flexible y gratuita. Sin embargo, su enfoque es el de una herramienta de dibujo general: permite representar visualmente un diagrama E/R, pero no interpreta semánticamente sus elementos, no valida reglas propias del modelo E/R y no genera una transformación relacional controlada.

Excalidraw es una herramienta de pizarra virtual colaborativa orientada a crear esquemas y diagramas con estética de dibujo manual~\cite{excalidraw}. Dispone de una versión gratuita y de una modalidad Excalidraw+ con funcionalidades adicionales~\cite{excalidrawPricing}. Puede resultar útil para bocetos o explicaciones informales, pero no está especializada en el modelo Entidad-Relación ni en la generación de esquemas relacionales o SQL.

Estas herramientas son útiles para dibujar o comunicar ideas, pero no cubren el objetivo principal de UBU E-R App: conectar la edición visual del diagrama con su validación, transformación relacional y generación SQL.

\section{Herramientas profesionales de modelado}

Lucidchart es una herramienta profesional de diagramación en línea que ofrece un plan gratuito y planes de pago orientados a usuarios individuales, equipos y organizaciones~\cite{lucidchartPricing}. Su ámbito es amplio y está pensado para distintos tipos de diagramas, colaboración y documentación visual.

Visual Paradigm ofrece herramientas específicas para modelado de bases de datos y creación de diagramas ERD. Su versión online incluye una herramienta gratuita para crear diagramas ERD, mientras que otras funcionalidades forman parte de planes o productos de pago~\cite{visualParadigmERD,visualParadigmPricing}.

Estas soluciones son más completas desde el punto de vista profesional, especialmente cuando se trabaja con modelos físicos, documentación empresarial o integración con otros procesos de diseño. No obstante, esa amplitud puede alejarlas del objetivo concreto de este TFG: proporcionar una herramienta docente, sencilla y centrada en la notación Entidad-Relación trabajada en la asignatura.

\section{ERDPlus}

ERDPlus es la alternativa más próxima al enfoque de UBU E-R App. Se trata de una herramienta web gratuita para crear diagramas Entidad-Relación, esquemas relacionales, esquemas en estrella y generar sentencias SQL~\cite{erdplus}. Además, su propia documentación presenta el servicio como una herramienta gratuita y durante la revisión manual se observó la presencia de publicidad en la interfaz~\cite{erdplus,erdplusTerms}.

Frente a herramientas de dibujo general, ERDPlus sí interpreta parte del contenido del diagrama y permite transformar modelos E/R en esquemas relacionales y SQL. Por este motivo, resulta especialmente relevante como punto de comparación. También incorpora soporte para elementos que no están completamente implementados en UBU E-R App, como ciertas formas de especialización o agregación.

No obstante, durante la revisión manual realizada para este trabajo se observaron varias diferencias importantes. Por ejemplo, algunas operaciones de edición resultan menos directas, la gestión de especializaciones puede ser rígida y la generación SQL simplifica decisiones que en un contexto docente son relevantes, como la distinción entre cardinalidad mínima cero y uno o el tratamiento de determinadas relaciones uno a uno mediante restricciones \texttt{UNIQUE}.

La Tabla~\ref{tab:comparativa_erdplus_ubu} resume las principales diferencias observadas entre ERDPlus y UBU E-R App desde el punto de vista del alcance de este trabajo.

\begin{table}[htbp]
    \centering
    \small
    \caption{Comparación entre ERDPlus y UBU E-R App.}
    \label{tab:comparativa_erdplus_ubu}
    \begin{tabular}{p{0.23\textwidth} p{0.31\textwidth} p{0.34\textwidth}}
        \toprule
        \textbf{Aspecto} & \textbf{ERDPlus} & \textbf{UBU E-R App} \\
        \midrule
        Disponibilidad
        & Herramienta web gratuita, con almacenamiento en la nube y publicidad.
        & Herramienta web gratuita, con despliegue web y posibilidad de ejecución local. \\

        Enfoque
        & Modelado E/R, esquemas relacionales, esquemas en estrella y generación SQL.
        & Modelado E/R académico, validación, transformación relacional y generación SQL simplificada. \\

        Elementos soportados
        & Soporta entidades, entidades débiles, atributos, cardinalidades y algunos elementos avanzados.
        & Soporta entidades fuertes y débiles, atributos compuestos y multivaluados, binarias, ternarias, reflexivas, roles, identificadoras e ISA inicial. \\

        ISA y agregaciones
        & Incluye soporte para especializaciones y mecanismos próximos a agregaciones o entidades asociativas.
        & ISA inicial con clave heredada; las agregaciones quedan fuera del alcance implementado. \\

        Transformación SQL
        & Genera SQL automáticamente, aunque en la revisión manual se observaron simplificaciones discutibles en algunos casos.
        & Genera SQL compatible con PostgreSQL, intentando respetar cardinalidades, claves primarias, claves foráneas, \texttt{UNIQUE}, \texttt{NOT NULL} y dependencias entre tablas. \\

        Usabilidad
        & Interfaz especializada, dependiente de conexión y con publicidad.
        & Interfaz adaptada al proyecto, sin publicidad, con persistencia local, importación/exportación JSON, selección múltiple, ajuste de vista y estructuras de ejemplo. \\
        \bottomrule
    \end{tabular}
\end{table}

\section{Comparativa final}

En conjunto, las alternativas revisadas permiten situar UBU E-R App en un espacio intermedio. No pretende competir con herramientas profesionales de modelado de bases de datos ni con editores gráficos de propósito general. Su aportación se encuentra en ofrecer una herramienta gratuita, accesible, sin publicidad, ejecutable localmente y orientada al aprendizaje del modelo Entidad-Relación, manteniendo conectadas la edición visual, la representación interna, la validación, la transformación relacional y la generación SQL.

La herramienta más cercana es ERDPlus, pero la comparación muestra que UBU E-R App resulta más ajustada al alcance docente de este TFG en aspectos como la ausencia de publicidad, la posibilidad de ejecución local, el soporte de relaciones ternarias con roles, la validación orientada al usuario y el control explícito de ciertos casos de transformación relacional.